\documentclass[11pt]{article}
\usepackage[margin=1in]{geometry}
\usepackage{amsmath,amssymb,amsthm,mathtools}
\usepackage{graphicx}
\usepackage{microtype}
\usepackage[hidelinks]{hyperref}

\newtheorem{theorem}{Theorem}
\newtheorem{lemma}{Lemma}
\newcommand{\R}{\mathbb R}
\newcommand{\Vfour}{C_2\times C_2}
\newcommand{\dist}{\operatorname{dist}}

\title{A sharper distortion bound for the minimal faithful\\
representation of $C_2\times C_2$}
\author{Candidate partial result for arXiv:2506.04425}
\date{27 August 2026}

\begin{document}
\maketitle

\begin{center}
\fbox{\parbox{0.9\linewidth}{\textbf{Status.} Candidate partial result,
likely valid. This packet improves the upper bound for one distinguished
representation; it does not compute the group contortion.}}
\end{center}

\section{Source question and contribution}

Blum-Smith, Derksen, Mixon, Qaddura, and Vose define the Euclidean contortion
of a finite group by
\[
  \Upsilon(G)=\sup_{V\in\operatorname{Rep}(G)}c_2(V/G),
\]
and ask in Section 4.2 for the values of $\Upsilon(C_4)$ and
$\Upsilon(\Vfour)$ \cite{BDMQV}. The question appears on page 43:

\begin{center}
\includegraphics[width=\linewidth]{figures/open_problem_crop.png}
\end{center}

Here $c_2(X)$ is the least distortion of a bilipschitz embedding of the metric
space $X$ into a Hilbert space. The source proves
$\sqrt2\leq\Upsilon(\Vfour)\leq2$. We study the natural minimal faithful
representation
\[
G:=\{\operatorname{diag}(1,1,1),\operatorname{diag}(1,-1,-1),
\operatorname{diag}(-1,1,-1),\operatorname{diag}(-1,-1,1)\}\curvearrowright\R^3.
\]
This is the multiplicity-free direct sum of the three nontrivial real
characters of $\Vfour$.

\section{Proof intuition}

An orbit is determined by the three absolute coordinates and, away from the
coordinate planes, the parity of the signs. Thus $\R^3/G$ is two copies of the
positive orthant glued along their boundary. On the same sheet the quotient
distance is Euclidean. Between sheets, the best group element leaves exactly
one sign mismatch, in the coordinate minimizing $a_i b_i$.

This suggests embedding the glued orthants by keeping $a=(|x_1|,|x_2|,|x_3|)$
and adding one signed coordinate equal to
$q(a):=\min_i a_i$, the Euclidean distance from $a$ to the boundary. The
resulting same-sheet and cross-sheet squared distances are explicit. A
two-variable scalar inequality bounds their ratio between $2-\sqrt2$ and $2$.
The quotient of these two constants is $2+\sqrt2$.

The original contortion problem remains open because a general
$\Vfour$-representation may contain arbitrary multiplicities of all three
nontrivial characters.

\section{Metric model}

Put $C=\R_+^3$ and $q(a)=\min_i a_i$. For $a\in C$ with all coordinates
positive, let $\varepsilon\in\{\pm1\}$ record the sign of the product of the
coordinates of a representative. If $q(a)=0$, this sign is immaterial because
the orbit contains representatives of both parities. Hence the orbit space is
identified with
\[
 X=(C\times\{\pm1\})/\big((a,+1)\sim(a,-1)\text{ when }q(a)=0\big).
\]

\begin{lemma}[Exact quotient metric]\label{lem:metric}
For $a,b\in C$, set
\[
 A=\|a-b\|_2^2,\qquad m=\min_{1\leq i\leq3}a_i b_i.
\]
Then
\[
d_X((a,\varepsilon),(b,\delta))^2=
\begin{cases}
A,&\varepsilon=\delta,\\
A+4m,&\varepsilon=-\delta.
\end{cases}
\]
\end{lemma}

\begin{proof}
For equal parity, an even sign change aligns all coordinate signs, giving
distance $\|a-b\|_2$, and coordinatewise absolute values show that no smaller
distance is possible. For opposite parity, any alignment has an odd number of
sign mismatches. If the mismatched coordinates form $S$, its squared distance
is
\[
A+4\sum_{i\in S}a_i b_i.
\]
All summands are nonnegative, so the minimum over odd $S$ is obtained by a
singleton at an index attaining $m$.
\end{proof}

\section{The embedding and its bound}

Define
\[
 F:X\longrightarrow\R^4,\qquad F(a,\varepsilon)=(a,\varepsilon q(a)).
\]
It is well-defined on the glued boundary since $q(a)=0$ there.

\begin{lemma}[The scalar inequality]\label{lem:scalar}
For $a,b\in C$, with $A,m$ as in Lemma~\ref{lem:metric}, and
$x=q(a),y=q(b)$, one has
\[
 A+(1+\sqrt2)(x+y)^2\geq4\sqrt2\,m.
\]
\end{lemma}

\begin{proof}
The assertion is immediate if $m=0$. By homogeneity, assume $m=1$. Then
$a_i b_i\geq1$ for each $i$, and $xy\leq1$. Define
\[
 f(t)=\begin{cases}(t^{-1}-t)^2,&0<t<1,\\0,&t\geq1.\end{cases}
\]
Choose an index $r$ with $a_r=x$. If $x<1$, then $b_r\geq1/x$, so
$(a_r-b_r)^2\geq f(x)$. Similarly, an index where $b$ attains $y$ contributes
at least $f(y)$. If the two indices coincide, $xy=1$, and at most one of
$f(x),f(y)$ is nonzero. Consequently,
\[
A\geq f(x)+f(y).\tag{1}
\]

Write $\lambda=1+\sqrt2$. First suppose $x,y<1$ and put $s=xy\leq1$.
Using $x^{-2}+y^{-2}\geq2/s$, $x^2+y^2\geq2s$, and
$(x+y)^2\geq4s$, (1) gives
\begin{align*}
A+\lambda(x+y)^2
&\geq \frac2s-4+(2+4\lambda)s\\
&\geq4\sqrt2.
\end{align*}
Indeed, the last one-variable expression is minimized at
$s=(1+\sqrt2)^{-1}$.

If, say, $x\geq1$, then $y\leq1$. Decreasing $x$ to $1$ only decreases the
remaining lower bound, and
\[
f(y)=\left(\frac1y-y\right)^2\geq4(1-y)^2.
\]
Thus
\[
A+\lambda(x+y)^2
\geq4(1-y)^2+\lambda(1+y)^2
\geq\frac{16\lambda}{4+\lambda}\geq4\sqrt2.
\]
The case $y\geq1$ is symmetric; if both are at least $1$, then $x=y=1$ and
the conclusion is immediate.
\end{proof}

\begin{theorem}[Improved bound for the minimal faithful representation]
\label{thm:main}
For the above action of $G\cong\Vfour$ on $\R^3$,
\[
 c_2(\R^3/G)\leq\sqrt{2+\sqrt2}.
\]
Together with the lower bound in Theorem 24 of \cite{BDMQV}, this gives
\[
 \sqrt2\leq c_2(\R^3/G)\leq\sqrt{2+\sqrt2}<2.
\]
\end{theorem}

\begin{proof}
For points on the same sheet, Lemma~\ref{lem:metric} and the fact that
$q=\min_i a_i$ is $1$-Lipschitz give
\[
d_X^2=A\leq\|F(a,\varepsilon)-F(b,\varepsilon)\|^2
=A+(x-y)^2\leq2A.
\]

For points on opposite sheets, put
\[
D^2=\|F(a,\varepsilon)-F(b,-\varepsilon)\|^2=A+(x+y)^2,
\qquad d_X^2=A+4m.
\]
Since $A\geq(x-y)^2$ and $m\geq xy$,
\[
2d_X^2-D^2=A+8m-(x+y)^2\geq4xy\geq0.
\]
Set $c=2-\sqrt2=\sqrt2(\sqrt2-1)$. Lemma~\ref{lem:scalar} yields
\begin{align*}
D^2-cd_X^2
&=(\sqrt2-1)A+(x+y)^2-4(2-\sqrt2)m\\
&=(\sqrt2-1)\big(A+(1+\sqrt2)(x+y)^2-4\sqrt2 m\big)\geq0.
\end{align*}
Thus every squared distance ratio lies in $[2-\sqrt2,2]$. Therefore
\[
\kappa(F)^2\leq\frac{2}{2-\sqrt2}=2+\sqrt2,
\]
as claimed.
\end{proof}

\section{Optimality inside the signed-boundary family}

For $t>0$, consider $F_t(a,\varepsilon)=(a,t\varepsilon q(a))$.
Same-sheet pairs that differ only in their smallest coordinate have squared
distance ratio exactly $1+t^2$. On the other hand, take opposite-sheet points
\[
a=(1,L,L),\qquad b=(L,1,L),\qquad L=1+\sqrt2.
\]
Here $A=4$, $m=L$, and $q(a)=q(b)=1$, so their squared distance ratio is
\[
\frac{4+4t^2}{4+4L}=\frac{1+t^2}{2+\sqrt2}.
\]
It follows that $\kappa(F_t)^2\geq2+\sqrt2$ for every $t>0$. Hence the choice
$t=1$ in Theorem~\ref{thm:main} is optimal throughout this natural
one-parameter family.

\section{Verification and limitations}

The script \texttt{code/verify\_bound.py} samples random pairs in both sheets,
checks the exact quotient formula by enumeration of the four group elements,
checks the proved ratio interval, and evaluates the extremal family above.
It is a regression check, not part of the proof.

The theorem does not determine $c_2(\R^3/G)$ exactly and does not upper-bound
$\Upsilon(\Vfour)$ below $2$, since contortion takes the supremum over all
finite-dimensional representations. The exact phrases ``Euclidean
contortion'', ``Upsilon(C4)'', and ``C2 times C2'', together with the paper
title and arXiv id, were searched on 27 August 2026. No later resolution or
matching bound was found. This is a bounded novelty check only.

\paragraph{Human review recommendation.}
Check Lemma~\ref{lem:scalar}, especially the case where the minima of $a$ and
$b$ occur at the same coordinate, and independently verify the cross-sheet
metric formula. If those pass, retain as a partial result; do not relabel it as
a solution of the contortion question.

\begin{thebibliography}{1}
\bibitem{BDMQV}
B. Blum-Smith, H. Derksen, D. G. Mixon, Y. Qaddura, and B. Vose,
\emph{Estimating the Euclidean distortion of an orbit space},
arXiv:2506.04425v1 (2025).
\end{thebibliography}

\end{document}
